%%%%%%%%%%%%%%%%%%%%%%%%%
\spellentry{Water Walk}
%%%%%%%%%%%%%%%%%%%%%%%%%

\tagline{Transmutation [Water]}
\begin{description*}
\item[Level:] Clr 3, Rgr 3
\item[Casting Time:] 1 standard action
\item[Range:] Touch
\item[Duration:] 10 min./level (D)
\item[Saving Throw:] Will negates (harmless)
\item[Spell Resistance:] Yes (harmless)
\item[Components:] V, S, DF
\end{description*}

The transmuted creatures can tread on any liquid as if it were firm ground. Mud,
oil, snow, quicksand, running water, ice, and even lava can be traversed easily,
since the subjects' feet hover an inch or two above the surface. (Creatures crossing
molten lava still take damage from the heat because they are near it.) The subjects
can walk, run, charge, or otherwise move across the surface as if it were normal
ground.
If the spell is cast underwater (or while the subjects are partially or wholly
submerged in whatever liquid they are in), the subjects are borne toward the surface
at 60 feet per round until they can stand on it.